% =====================================================================
% Chapter 2 -- Background and Related Work
% =====================================================================
\documentclass[../preamble.tex]{subfiles}
\begin{document}

\chapter{Background and Related Work}
\label{ch:background}

\section{Clinical aspects of PCOS}
\label{sec:bg-clinical}

Polycystic Ovary Syndrome is one of the most prevalent endocrine disorders
in women of reproductive age, with global frequencies ranging from
approximately 8\% to 12\% depending on the diagnostic criteria applied
\cite{ghaderzadeh2025,patil2025,mahesswari2024}. The Rotterdam consensus
requires two of three features after excluding alternative causes:
oligo- or anovulation, clinical or biochemical hyperandrogenism, and
polycystic ovarian morphology on ultrasound \cite{rotterdam2003}. The
international evidence-based guideline synthesises the diagnostic and
management recommendations used by clinicians \cite{teede2018}.

The clinical features and ultrasound morphology often co-vary with metabolic
and cardiovascular risk factors, so underdiagnosis and misdiagnosis carry
consequences that extend beyond the immediate ultrasound finding
\cite{ghaderzadeh2025,nsugbe2023}. A prediction tool that flags PCOS
suspicion can serve as a triage step, but it cannot substitute for a clinical
workup.

\section{Image-based machine learning for PCOS}
\label{sec:bg-imaging-ml}

Early PCOS machine-learning work relied on clinical features extracted from
questionnaires, blood panels, and limited imaging variables. Examples include
XGBoost trained on combined clinical and ultrasound features
\cite{agirsoy2025}, hybrid feature selection with ensemble classifiers
\cite{patil2025}, two-level random forests with Shapash explanations
\cite{mahesswari2024}, staged probabilistic inference with SMOTE
\cite{nsugbe2023}, and radial pulse-wave analyses for PCOS screening
\cite{lim2023}. These studies report high accuracy, but several are limited to
a single source and rely on tabular features rather than ultrasound
appearance.

A growing body of work focuses directly on ovarian ultrasound. Sundari
et~al.\ report 94.8\% accuracy, 93.2\% sensitivity, and 95.5\% specificity on
the PCOSgen cohort using an attention-enhanced EfficientNet-B3
\cite{sundari2025}. CystNet uses watershed segmentation and a convolutional
autoencoder for follicle-level features and reports 96.54--97.75\% accuracy
\cite{moral2024}. Lakshmi and Pushpa fuse ultrasound with clinical features
through a cross-attention module \cite{lakshmi2026}. Shanmugavadivel
et~al.\ combine a clinical branch with a CNN imaging branch
\cite{shanmugavadivel2024}. Reka et~al.\ apply ESRGAN super-resolution and
the Segment Anything Model before classification \cite{reka2025}.
Jeyashanker et~al.\ introduce a custom CoCo-HoloNet with adaptive
optimisation \cite{jeyashanker2025}. Bedi et~al.\ pair adaptive bilateral
filtering with an attention residual U-Net \cite{bedi2024}. Fan
et~al.\ design a lightweight Ocys-Net architecture \cite{fan2023}. Ghosh and
Srinivasan ensemble EfficientNetB7 with DenseNet201 and use genetic-algorithm
hyperparameter optimisation \cite{ghosh2025}. Poorani and Khilar explore
multiscale thresholding as a non-deep segmentation alternative
\cite{poorani2024}.

These studies report strong internal metrics, but several do not evaluate on
a separate external cohort. A systematic review found that only about 25\% of
PCOS deep-learning studies applied any form of explainability and that the
overuse of a small set of public datasets limits generalisability
\cite{ghaderzadeh2025}. Recent systematic surveys of the broader
PCOS machine-learning literature reach similar conclusions
\cite{suha2023,ghosh2025}.

\section{Calibration and uncertainty}
\label{sec:bg-calib-unc}

Modern neural networks are systematically overconfident on classification
tasks \cite{guo2017}. Among the standard post-hoc calibration methods,
temperature scaling, which learns a single scalar multiplier on the logits,
performs close to the strongest alternatives while keeping the original
accuracy. The thesis uses temperature scaling both per-model and on the
averaged ensemble probability, an approach consistent with the broader
calibration literature \cite{guo2017}. Bayesian extensions of CNNs provide a
separate, complementary view of uncertainty through MC-Dropout
\cite{gal2016}. Akram et~al.\ compare MC-Dropout with mean-field variational
inference and deterministic inference for ultrasound classification
\cite{akram2025}, while Roy et~al.\ demonstrate voxel-level uncertainty maps
for brain MRI segmentation \cite{roy2018}. The thesis uses MC-Dropout
because it requires no architectural changes, integrates with the existing
training loop, and has been widely adopted in medical-imaging uncertainty
work.

\section{Explainability}
\label{sec:bg-xai}

Gradient-based explanations such as Grad-CAM highlight image regions that
influenced a prediction and can reveal shortcut behaviour
\cite{selvaraju2017}. Excitation backprop and layer-wise relevance
propagation offer alternative signal propagation schemes
\cite{zhang2018,samek2017}. Tjoa and Guan survey explainability methods and
the gap between technical interpretability and clinical interpretability
\cite{tjoa2021}. The executed pipeline uses Grad-CAM only; LRP and SHAP were
proposed in the original FYDP brief but were not executed in the final
analysis. This is recorded in \cref{app:scope-evolution}.

\section{Preprocessing for ultrasound}
\label{sec:bg-preprocessing}

Ultrasound images contain multiplicative speckle noise and low-contrast
structures. Anisotropic diffusion methods were developed to smooth
homogeneous regions while preserving edges. SRAD replaces the
gradient-magnitude diffusion coefficient used in Perona--Malik with the
instantaneous coefficient of variation, which is appropriate for
multiplicative noise \cite{yu2002}. CLAHE is a standard contrast-enhancement
step for low-contrast medical images. This thesis combines a fixed CLAHE
step with a controlled SRAD-versus-Gaussian denoising comparison, so the
contribution of the diffusion choice can be attributed without confounding
from other preprocessing steps.

\section{Summary of literature gaps}
\label{sec:bg-gaps}

Three gaps motivate this work:

\begin{enumerate}
  \item Most PCOS deep-learning studies evaluate on a single source, so it is
        unclear how the reported metrics would transfer to a different cohort.
  \item Few studies report calibration, MC-Dropout, or any principled
        uncertainty metric on PCOS ultrasound.
  \item Few studies combine cross-dataset evaluation with explainability on
        the same cohort using the same model checkpoints.
\end{enumerate}

The thesis addresses these gaps by combining external validation, two-pass
calibration, MC-Dropout entropy, and Grad-CAM analysis within a single
reproducible pipeline.

\end{document}